% !TeX root = ../sections/02_exercises.tex
\subsection{1-A-1-1.}

\begin{exercise}
	% 問題文
	$0<r<1$に対して、$\lim_{n\rightarrow\infty}r^n=0$であることを\\$\epsilon N$ 論法を用いて証明せよ.

\end{exercise}

\begin{background}
	この問題で用いる基本事項を整理する。

	\begin{definition}[数列の極限]
		数列 $\{a_n\}$ が実数 $\alpha$ に収束するとは，
		任意の $\epsilon>0$ に対して，ある自然数 $N$ が存在して，
		すべての $n\geq N$ に対して
		\[
			|a_n-\alpha|<\epsilon
		\]
		が成り立つことをいう。
		このとき，
		\[
			\lim_{n\to\infty}a_n=\alpha
		\]
		と書く。
	\end{definition}

	\begin{remark}
		今回の問題では
		\[
			a_n=r^n,\qquad \alpha=0
		\]
		である。したがって，示すべき不等式は
		\[
			|r^n-0|<\epsilon
		\]
		である。
	\end{remark}

	\begin{remark}
		仮定より $0<r<1$ であるから，任意の自然数 $n$ に対して
		\[
			r^n>0
		\]
		である。よって
		\[
			|r^n-0|=r^n
		\]
		となる。
		したがって，この問題では
		\[
			r^n<\epsilon
		\]
		を示せばよい。
	\end{remark}

	\begin{theorem}[アルキメデスの性質]
		任意の実数 $A$ に対して，ある自然数 $N$ が存在して
		\[
			N>A
		\]
		が成り立つ。
	\end{theorem}

	\begin{lemma}
		$0<r<1$ のとき，
		\[
			\log r<0
		\]
		である。
	\end{lemma}

	\begin{remark}
		$\log r<0$ であるため，不等式の両辺を $\log r$ で割ると，
		不等号の向きが反転することに注意する。
	\end{remark}
\end{background}

\begin{strategy}
	まず，極限の定義に従って考える。

	示したいことは
	\[
		\lim_{n\to\infty}r^n=0
	\]
	である。したがって，任意の $\epsilon>0$ に対して，
	ある自然数 $N$ が存在し，すべての $n\geq N$ に対して
	\[
		|r^n-0|<\epsilon
	\]
	が成り立つことを示せばよい。

	ここで $0<r<1$ より $r^n>0$ であるから，
	\[
		|r^n-0|=r^n
	\]
	である。したがって目標は
	\[
		r^n<\epsilon
	\]
	を示すことに変わる。

	次に，この不等式を満たすような $n$ の条件を逆算する。
	両辺の対数をとると，
	\[
		r^n<\epsilon
	\]
	は
	\[
		n\log r<\log\epsilon
	\]
	となる。

	ただし $0<r<1$ より
	\[
		\log r<0
	\]
	である。したがって，両辺を $\log r$ で割ると不等号の向きが反転し，
	\[
		n>\frac{\log\epsilon}{\log r}
	\]
	を得る。

	よって，アルキメデスの性質より，
	\[
		N>\frac{\log\epsilon}{\log r}
	\]
	を満たす自然数 $N$ を取ることができる。

	このように選べば，すべての $n\geq N$ に対して
	\[
		n>\frac{\log\epsilon}{\log r}
	\]
	となり，逆にたどることで
	\[
		r^n<\epsilon
	\]
	が従う。したがって，
	\[
		|r^n-0|<\epsilon
	\]
	が成り立つ。
\end{strategy}